% Version 0.25.  The actual algebraic K-theory bridge and its quantitative join.
\section{From algebraic \texorpdfstring{$K$}{K}-theory to the uniform second-cohomology bound}
\label{sec:k-theory-bridge}

\status{KB1--KB4 construct the literal \texttt{KBridge} using classical
algebraic $K$-theory on the given number field.  KB5--KB8 close its
cardinality and constant chain: together with LD3, LK2--LK3 and TB10,
they prove the literal statement \texttt{LemmaK 31}.  Quillen finite
generation, Borel's rank theorem, Dwyer--Friedlander surjectivity,
Arlettaz's Hurewicz bound and integral homology stability are explicit
classical citation leaves with their hypotheses matched below.  None
is inferred merely from the name of a \texttt{Published} interface.}

Fix a prime $\ell$, a totally complex number field $L$, its specified
algebraic closure $\Omega$, and an arbitrary datum
$D:\texttt{LData}\ L\ \Omega\ \ell$.  Thus $S=D.S$ is exactly the set
of finite places above $\ell$, $D.A_v$ is the chosen valuation subring
above each $v\in S$, and $D.\chi:G_{L,S}\to\mathbb Z_\ell^\times$
pulls back to the cyclotomic character on $G_L$.
Assume \texttt{HasMuQ}: $L$ contains a primitive $\ell$-th root if
$\ell$ is odd, and a primitive fourth root if $\ell=2$.
Throughout put
\[
\begin{gathered}
 n=[L:\mathbb Q],\quad \Delta=|\operatorname{disc}L|,\quad
 e\ge1,\quad m=2e,\quad N=4e+4=2m+4,\\
 X=H^2_{\rm cts}(G_{L,S},\mathbb Z_\ell(e+1)).
\end{gathered}
 \tag{KB.1}
\]
The cohomology in (KB.1) uses continuous inhomogeneous cochains and
the compact $\ell$-adic topology on $\mathbb Z_\ell$.
Write $h_L$ for the ordinary ideal class number, and
\[
 H=H_m(\operatorname{GL}_N(\mathcal O_L),\mathbb Z),\qquad
 T=H_{\rm tors},\qquad
 c(m)=\sum_{a=1}^{m}(a+2)=\frac{m(m+5)}2.
 \tag{KB.2}
\]
These are respectively the source's \texttt{HGL}, \texttt{torsHGL}
and \texttt{cOne}; all group homology has trivial integral coefficients.

\subsection{Matching the original cochains and the arithmetic comparison}

\begin{proposition}[KB1: original datum and integral cochain comparison]
\label{claim:KB1}
Let $U=\operatorname{Spec}\mathcal O_L[1/\ell]$ and let $j$ be an
integer.  With the character supplied by $D$, there are compatible
identifications
\[
 H^q_{\rm et}(U,\mathbb Z/\ell^r(j))
       \simeq H^q_{\rm cts}(G_{L,S},\mathbb Z/\ell^r(j)),
 \qquad
 H^2_{\rm cts}(G_{L,S},\mathbb Z_\ell(j))
       \simeq\varprojlim_r H^2_{\rm et}(U,\mathbb Z/\ell^r(j)).
 \tag{KB.3}
\]
The finite-level groups are finite; they vanish for $q\ge3$.
Moreover, directly on these cochains,
\[
 H^2_{\rm cts}(G_{L,S},\mathbb Z_\ell(j))/\ell
       \lhook\joinrel\longrightarrow
 H^2_{\rm cts}(G_{L,S},\mathbb F_\ell(j)).
 \tag{KB.4}
\]
No vanishing assertion about integral $H^3$ is needed for (KB.4).
\end{proposition}

\begin{proof}
The source's $G_{L,S}$ is the quotient of $G_L$ by the closed normal
subgroup generated by inertia at the omitted finite places.  Its
finite quotients are therefore exactly the finite Galois extensions
unramified outside $S$.  This identifies it with $\pi_1(U)$, using
the geometric point specified by $\Omega$.  Since $G_L\to G_{L,S}$
is onto, the identity $D.\chi\circ(G_L\to G_{L,S})=\chi_\ell$
determines $D.\chi$ uniquely.  It is continuous by the quotient
topology.  Consequently the action in \texttt{zlTw D j} is precisely
$\chi_\ell^j$, including for negative $j$, and reduction gives the
usual finite Tate twist.  Restriction to each original $D.A_v$
commutes with these identifications.

The finite comparison is Milne, \emph{Arithmetic Duality Theorems},
II, Proposition~2.9, applied to the locally constant finite sheaf
$\mathbb Z/\ell^r(j)$; $\ell$ is invertible on $U$.
Finiteness is II, Theorem~2.13.  The higher-degree vanishing is the
finite global conclusion already matched in LD3, since $L$ has no
real places.\footnote{J.S. Milne, \emph{Arithmetic Duality Theorems},
2nd edition (2006), I, Theorem~4.10, and II, Proposition~2.9 and
Theorem~2.13; \url{https://jmilne.org/math/Books/ADTnot.pdf}.}

Here is the inverse-limit argument, to fix which integral cohomology
is meant.  Put $C_r^q=C_{\rm cts}(G_{L,S}^q,\mathbb Z/\ell^r(j))$.
The transition $C_{r+1}^q\to C_r^q$ is onto: any cochain has finitely
many values and can be lifted by choosing one residue lift for each
value.  The resulting lift is locally constant.  Also
\[
 C_{\rm cts}(G_{L,S}^q,\mathbb Z_\ell(j))=\varprojlim_r C_r^q;
 \tag{KB.5}
\]
a map into the inverse-limit topology is continuous exactly when
all its coordinate maps are continuous.  The short exact sequence
of complexes with product map
$(c_r)_r\mapsto(c_r-\operatorname{red}(c_{r+1}))_r$ gives
\[
 0\longrightarrow\varprojlim{}^1 H^{q-1}(C_r)
 \longrightarrow H^q(\varprojlim C_r)
 \longrightarrow\varprojlim H^q(C_r)\longrightarrow0.
 \tag{KB.6}
\]
For an inverse system of finite groups the product difference map
is onto: any finite set of its equations can be solved backwards
from its largest index, and compactness of the product of the finite
groups then solves all equations.  Thus its first derived limit
vanishes.  Apply this to the finite $H^1(C_r)$ to obtain (KB.3).

Finally the coefficient sequence
$0\to\mathbb Z_\ell(j)\xrightarrow{\ell}\mathbb Z_\ell(j)
\to\mathbb F_\ell(j)\to0$ is exact on continuous cochains.
Surjectivity uses the same finite-valued lift; division by $\ell$
on $\ell\mathbb Z_\ell$ is continuous, so exactness also holds in
the middle.  The cohomology exact sequence gives (KB.4).
In degree two its differential is
$dc(g,h)=g c(h)-c(gh)+c(g)$, exactly the source's differential and
continuous-coboundary quotient.  This matches the actual $X$, not
just an abstract group of the same order.
\end{proof}

\begin{proposition}[KB2: finite arithmetic K-group mapping onto the actual target]
\label{claim:KB2}
Let $A_m=K_m(\mathcal O_L)$ be the usual Quillen group, and set
$K_\ell=\mathbb Z_\ell\otimes_{\mathbb Z}A_m$.  Then $A_m$ and
$K_\ell$ are finite, and there is a surjective
$\mathbb Z_\ell$-linear map
\[
                    p:K_\ell\twoheadrightarrow X.
 \tag{KB.7}
\]
This includes $e=1$ and $\ell=2$.
\end{proposition}

\begin{proof}
Quillen's finite-generation theorem says that every
$K_i(\mathcal O_L)$ is finitely generated.  Borel's rank formula
gives rank zero for every positive even index.  Thus $A_m$ is a
finite abelian group, and $K_\ell$ is its $\ell$-primary part.
The groups here can be defined by
$K_i(\mathcal O_L)=\pi_i(B\operatorname{GL}_\infty(\mathcal O_L)^+)$
for $i\ge1$, the standard plus-construction model of Quillen
$K$-theory.\footnote{D. Quillen, ``Finite generation of the groups
$K_i$ of rings of algebraic integers'', \emph{Lecture Notes in
Mathematics} 341 (1973), 179--198, Theorem~1,
\url{https://doi.org/10.1007/BFb0067056}; original conference volume:
\url{https://www.sas.rochester.edu/mth/sites/doug-ravenel/otherpapers/bass-seattle.pdf}.
A. Borel, ``Stable real cohomology of arithmetic groups'',
\emph{Ann. Sci. ENS} 7 (1974), 235--272, Proposition~12.2,
\url{https://www.numdam.org/item/10.24033/asens.1269.pdf}.
The canonical plus model and its agreement with the exact-category
model are part of Quillen's $K$-theory; see also Arlettaz, cited in
KB3, Definition~3.12 and Remark~3.16.}

The arithmetic input is Dwyer--Friedlander's comparison.
Their Proposition~8.1 removes the primes over $\ell$ in finite
coefficients; Theorem~8.5 gives a surjection from algebraic to
\'{e}tale $K$-theory in degree $m\ge2$.  Their Proposition~8.2
also supplies the degree-two case directly.
For $\ell=2$ use the cofinal levels $r\ge2$ and the primitive
fourth root required by that theorem.  Our root is integral and
its square is $-1$.  Totally complex number fields have
$\ell$-cohomological dimension at most two, as required.
Their Proposition~5.2 supplies the finite-coefficient
\'{e}tale spectral sequence.\footnote{W.G. Dwyer and E.M. Friedlander,
``Algebraic and \'{e}tale $K$-theory'', \emph{Trans. Amer. Math. Soc.}
292 (1985), 247--280, Propositions~5.2, 8.1--8.2 and Theorem~8.5;
the completed conclusion is Theorem~8.7 and Remark~8.8.
\url{https://doi.org/10.1090/S0002-9947-1985-0805962-2}.
Author-uploaded full text:
\url{https://www.researchgate.net/publication/239378605_Algebraic_and_Etale_K-Theory}.}

We spell out the passage to compact coefficients.  The coefficient
exact sequence is
\[
 0\longrightarrow A_m/\ell^r A_m
 \longrightarrow K_m(\mathcal O_L;\mathbb Z/\ell^r)
 \longrightarrow K_{m-1}(\mathcal O_L)[\ell^r]
 \longrightarrow0.
 \tag{KB.8}
\]
The transition on its rightmost terms is multiplication by $\ell$,
as seen from the two-term free resolutions of $\mathbb Z/\ell^r$.
Since $K_{m-1}(\mathcal O_L)$ is finitely generated, its torsion has
bounded exponent, and that inverse limit is zero.  The leftmost
groups are finite and have vanishing first derived limit.  Hence
the inverse limit of the middle terms is $K_\ell$.

Write $E_r=K_m^{\rm et}(U;\mathbb Z/\ell^r)$.  The spectral sequence
has only cohomological columns $0,1,2$, by KB1, and only even
coefficient degrees.  In total degree $m=2e$ it therefore gives
\[
 0\longrightarrow H^2_{\rm et}(U,\mathbb Z/\ell^r(e+1))
 \longrightarrow E_r
 \longrightarrow H^0_{\rm et}(U,\mathbb Z/\ell^r(e))
 \longrightarrow0.
 \tag{KB.9}
\]
All these groups are finite.  The inverse limit of the last terms
is $H^0(G_{L,S},\mathbb Z_\ell(e))=0$.  Indeed the cyclotomic
image is infinite, since $[\mathbb Q(\mu_{\ell^r}):\mathbb Q]$
is unbounded whereas $[L:\mathbb Q]=n$ is fixed.  It contains a
unit $u$ with $u^e\ne1$; the torsion-free module
$\mathbb Z_\ell$ has no nonzero invariant under that unit.
By (KB.9) and KB1, $\varprojlim E_r\simeq X$.

The finite-level comparison maps in the preceding paragraph are
surjective and compatible.  A compatible target sequence has a
nonempty finite set of lifts at each level.  Every finite string
of compatibility conditions can be satisfied by lifting at its
largest level; compactness gives a compatible infinite string.
Their inverse-limit map is thus onto.  Combining it with (KB.8),
(KB.9) and (KB.3) constructs (KB.7).  All maps respect the scalar
actions, so the limit map is $\mathbb Z_\ell$-linear.
The cohomological-dimension assertion used for the fraction field
also follows from LD3: a finite-valued continuous cocycle on $G_L$
factors through an extension ramified at only finitely many places;
enlarge $S$ to contain those places and apply the higher-degree
vanishing there.  This uses no integral second-cohomology bound.
\end{proof}

\subsection{The homology target and the Hurewicz loss}

\begin{proposition}[KB3: integral stable Hurewicz map with the required kernel exponent]
\label{claim:KB3}
There is a $\mathbb Z_\ell$-linear map
\[
 h:K_\ell\longrightarrow\mathbb Z_\ell\otimes_{\mathbb Z}H
 \quad\hbox{such that}\quad
 h(x)=0\ \Longrightarrow\ \ell^{c(m)}x=0.
 \tag{KB.10}
\]
It is obtained from the canonical Hurewicz map and the actual block
stabilizations of the source.
\end{proposition}

\begin{proof}
We first check the ring hypothesis for stability.  For a Dedekind
domain $R$, let $(a_0,\ldots,a_t)$ be unimodular with $t\ge2$.
If all first $t$ entries vanish, $a_t$ is a unit, and adding it to
one of those entries shortens the row.  Otherwise choose $j<t$
with $a_j\ne0$, and choose $k<t$ different from $j$.  Only finitely
many maximal ideals contain $a_j$.  The Chinese remainder theorem
gives $b\in R$ which is $1$ modulo each such ideal containing
$a_k$, and $0$ modulo the remaining ones.  Replace $a_k$ by
$a_k+ba_t$, leaving the other first $t$ entries fixed.  A maximal
ideal containing the resulting row contains $a_j$.  If it does
not contain $a_k$, the choice $b=0$ is contradictory.  If it does,
the choice $b=1$ forces $a_t$ into it, contradicting the original
unimodularity.  This proves precisely \texttt{StableRangeLe R 2}.

Apply this to $R=\mathcal O_L$.  Van der Kallen's integral stability
theorem applies to block inclusions $\operatorname{GL}_a(R)
\hookrightarrow\operatorname{GL}_{a+1}(R)$.  In his Convention~2.2,
our stable range permits $\mathrm{sdim}=1$; his Theorem~4.11 gives
isomorphisms in degree $m$ for $a\ge2m+1$.
In particular all the maps starting at $N=2m+4$ are
isomorphisms.\footnote{W. van der Kallen, ``Homology stability for
linear groups'', \emph{Invent. Math.} 60 (1980), 269--295,
Convention~2.2 and Theorem~4.11 (p.~289).
The theorem states surjectivity for $a\ge2m+\max(1,\mathrm{sdim})-1$
and injectivity for $a\ge2m+\max(1,\mathrm{sdim})$.
Author's copy: \url{https://webspace.science.uu.nl/~kalle101/homstab1980.pdf}.}
The bar complex commutes with the union of these groups: each bar
chain involves finitely many matrices, all in some finite rank.
Filtered unions are exact for abelian groups, so
\[
 H\xrightarrow{\ \sim\ }
 H_m(\operatorname{GL}_\infty(\mathcal O_L),\mathbb Z).
 \tag{KB.11}
\]
The block maps are exactly \texttt{stab}, and the union is the
source's \texttt{GLinf}.  Singular homology of the nerve model
$B\operatorname{GL}_\infty$ agrees with this bar homology.

The canonical plus construction preserves integral homology.
Compose its degree-$m$ Hurewicz map with that identification to get
\[
 f:A_m\longrightarrow
 H_m(\operatorname{GL}_\infty(\mathcal O_L),\mathbb Z).
\]
Arlettaz's Corollary~7.19(a) states that $R_m\ker f=0$, with
$R_m=\prod_{a=2}^m L_a$ and every $L_a$ squarefree.\footnote{
D. Arlettaz, ``Algebraic $K$-theory of rings from a topological
viewpoint'', \emph{Publicacions Matem\`atiques} 44 (2000), 3--84,
Definitions~7.11 and 7.14 and Corollary~7.19(a), pp.~52--55.
\url{https://mat.uab.es/pubmat/fitxers/download/FileType:pdf/FolderName:v44(1)/FileName:44100_01.pdf}.}
Thus
\[
                   v_\ell(R_m)\le m-1\le c(m).
 \tag{KB.12}
\]
For an exact check of the source constant, $L_a$ is the product of
primes $p$ admitting a nonnegative integer sequence $(b_i)_{i\ge0}$
with sum $a$, $b_0\equiv0\pmod{2p-2}$,
$b_i\equiv0\text{ or }1\pmod{2p-2}$ for $i>0$, and
$p b_{i+1}\le b_i$.  These are \texttt{ArlAdm} and \texttt{arlL}.
The sequence is finitely supported because its sum is finite;
for $a>0$ one has $b_0\ge2p-2$, so the source's bound $p\le a$
omits no prime.  Also $L_1=1$, so its \texttt{arlR m} is this
same $R_m$.  Only squarefreeness of each factor is used in (KB.12).

The torsion-free abelian group $\mathbb Z_\ell$ is flat over
$\mathbb Z$.  Tensoring the kernel sequence for $f$ therefore shows
that $R_m$ annihilates the kernel of $1\otimes f$.  Write
$R_m=\ell^a u$, with $a=v_\ell(R_m)$ and $u$ an $\ell$-adic unit.
Multiplication by $u$ is invertible, so that kernel is killed by
$\ell^a$, hence by $\ell^{c(m)}$.  Compose $1\otimes f$ with
the tensor of the inverse of (KB.11).  This is (KB.10), including
all small primes; no prime is discarded from the homology group.
\end{proof}

\begin{proposition}[KB4: the actual KBridge supplier]
\label{claim:KB4}
The literal proposition \texttt{KBridge} holds.
For every original $\ell,L,\Omega,D,e$ satisfying its hypotheses,
one may take the finite $\mathbb Z_\ell$-module $K_\ell$ of KB2,
the surjection $p$ of (KB.7), and the map $h$ of (KB.10).
In particular $X$ is finite.
\end{proposition}

\begin{proof}
KB2 supplies the finite module and surjectivity.  KB3 supplies the
map (KB.10), with the specific homology group $H$ of (KB.2),
and the exponent $c(2e)$.
These are all the fields in \texttt{KBridge}.  Every coefficient
identification in KB1 used the original $D$ and $\Omega$, and no
new local places were selected.  Finiteness of $X$ follows by
surjectivity from the finite $K_\ell$.

The source files \texttt{KBridgeOfKPublished} and
\texttt{KBridgeOfKPublishedW} specify packages of interfaces.
This proof constructs the existential bridge itself with canonical
$K$-theory.  It therefore does not need, or assert, a universal
comparison theorem for every arbitrary record satisfying the
source's definition of a \texttt{Plus} model.
\end{proof}

\subsection{Finite groups, strict local kernels, and the class number}

\begin{proposition}[KB5: the finite-module and tensor cardinality calculation]
\label{claim:KB5}
Suppose $K$ is a finite $\mathbb Z_\ell$-module,
$p:K\twoheadrightarrow X$ and $h:K\to W$ are linear, and
$\ell^c\ker h=0$.  Then
\[
                  |X|\le |X/\ell X|^c\,|\operatorname{im}h|.
 \tag{KB.13}
\]
If $W=\mathbb Z_\ell\otimes_{\mathbb Z}H$ and $T=H_{\rm tors}$
is finite, every finite $\mathbb Z_\ell$-submodule $Y\subset W$
satisfies $|Y|\le|T|$.  Finite generation of $H$ is not required.
\end{proposition}

\begin{proof}
Put $Y=p(\ker h)$.  Then $\ell^cY=0$ and $X/Y$ is a quotient
of $K/\ker h\simeq\operatorname{im}h$.  For a finite group $X$,
multiplication by $\ell^c$ has kernel and cokernel of equal order.
Moreover multiplication by $\ell^i$ maps $X/\ell X$ onto
$\ell^iX/\ell^{i+1}X$.  The filtration of $X/\ell^cX$ therefore
gives
\[
 |Y|\le|X[\ell^c]|=|X/\ell^cX|
                \le|X/\ell X|^c.
\]
Multiplication by $|X/Y|$ proves (KB.13), also when $c=0$.

For the tensor assertion, $H/T$ is torsion-free and hence flat over
$\mathbb Z$.  Tensor its flatness with the injection
$\mathbb Z_\ell\hookrightarrow\mathbb Q_\ell$ to see that
$\mathbb Z_\ell\otimes(H/T)$ embeds in
$\mathbb Q_\ell\otimes(H/T)$; it is thus torsion-free as an
abelian group.  Flatness of $\mathbb Z_\ell$ gives the exact sequence
\[
 0\longrightarrow\mathbb Z_\ell\otimes T
 \longrightarrow\mathbb Z_\ell\otimes H
 \longrightarrow\mathbb Z_\ell\otimes(H/T)\longrightarrow0.
\]
The image of finite $Y$ in the last group is zero.
Decomposing the finite abelian group $T$ into cyclic groups shows
that $\mathbb Z_\ell\otimes T$ is its $\ell$-primary part:
$\mathbb Z_\ell\otimes\mathbb Z/a\mathbb Z
\simeq\mathbb Z_\ell/a\mathbb Z_\ell$.
Its order is at most $|T|$, which proves the claim.
\end{proof}

\begin{proposition}[KB6: the generator bound on the same arithmetic datum]
\label{claim:KB6}
For every integer $j$, with the datum and roots of unity fixed above,
\[
 \left|H^2_{\rm cts}(G_{L,S},\mathbb Z_\ell(j))/\ell\right|
                         \le h_L\ell^n.
 \tag{KB.14}
\]
Here finiteness of the left side follows from KB1, independently of
the integral finiteness supplied by KB4.
\end{proposition}

\begin{proof}
Let $M=\mathbb F_\ell(j)$ and
$M^D=\operatorname{Hom}_{\mathbb F_\ell}(M,\mathbb F_\ell(1))$.
The dual action is $\chi_\ell(g)\varphi(g^{-1}x)$, so
$M^D=\mathbb F_\ell(1-j)$.  This is precisely the inverse-and-twist
convention in \texttt{zd}.  LD3 gives a perfect finite pairing
between the strict local kernels
$\operatorname{Sha}^2_S(M)$ and $\operatorname{Sha}^1_S(M^D)$;
their orders are equal.

We make the local-field join explicit.  Keep $A_v=D.A_v$ and let
$L_v^h\subset\Omega$ be the henselization selected by this
valuation, whose absolute Galois group is $\operatorname{Dec}(A_v)$.
Its completion is $L_v$.  Completion gives an equivalence of finite
separable extensions of $L_v^h$ and $L_v$, with the same embeddings
and degrees; this is the henselian completion comparison obtained
from Krasner's lemma.\footnote{This is the finite-extension
comparison already used in the setup of LD3.  For its precise
henselian hypotheses, see Milne, \emph{Arithmetic Duality Theorems},
I, Appendix A, Case~2, pp.~130--131: henselizations of local rings
of global fields are excellent, and completion preserves the
finite separable extensions and their Galois groups.}
Embed the finite $\mathbb Q_\ell$-extension $L_v$ into
$\texttt{PadicAlgCl}(\ell)$ and denote its image by $K_v$.
Transporting the separable closures gives a continuous group
isomorphism, with continuous inverse,
\[
 \operatorname{Dec}(A_v)\simeq
 \operatorname{Gal}\bigl(\texttt{PadicAlgCl}(\ell)/K_v\bigr).
\]
It preserves the action on every $\ell$-power root of unity;
therefore the restriction of the original cyclotomic character
becomes exactly \texttt{cycK}.  Precomposition in each cochain
degree commutes with the differential, cup product and the
original restriction $\operatorname{Dec}(A_v)\to G_{L,S}$.
This is the source's
\texttt{decompositionGroup\_mulEquiv\_localGalois\_cyclotomic}
join, on the chosen $A_v$.  LK1 identifies its local cochains,
and LK2--LK3, with $r=1$, consequently give
\[
 |H^2(L_v,M)|=|H^0(L_v,M^D)|\le\ell \quad(v\in S).
\]
The kernel of localization $H^2(G_{L,S},M)\to
\prod_{v\in S}H^2(L_v,M)$ is exactly
$\operatorname{Sha}^2_S(M)$.  Taking its image, which may be a
proper subgroup of the product, yields
\[
 |H^2(G_{L,S},M)|
       \le |\operatorname{Sha}^1_S(M^D)|\ell^{|S|}.
 \tag{KB.15}
\]
The sum of the products of ramification and residue degrees at
$\ell$ is $n$, and each product is at least one; hence $|S|\le n$.

The reduction of the cyclotomic character modulo $\ell$ is trivial:
all $\ell$-th roots already belong to $L$.  At $\ell=2$ this also
follows directly from $\mathbb F_2^\times=1$.
Thus $M^D$ is the trivial one-dimensional module.  Evaluation of a
dual linear functional at $1$ identifies it with $\mathbb F_\ell$
and commutes with all the original local restriction maps.
Its degree-one cocycles are continuous additive characters, and
its coboundaries are zero.  A strict local-kernel character is
zero on the chosen decomposition groups at every $v\mid\ell$.
Inflated to $G_L$, it is unramified outside $S$ by definition and
at $S$ by this zero restriction.  Conjugate decomposition groups
give the same conclusion, because its target is abelian.

Global unramified class field theory now factors this character
through $\operatorname{Cl}(\mathcal O_L)$.  There is no real-place
condition, since $L$ is totally complex.\footnote{The precise
classical input is the Artin identification of the maximal abelian
extension unramified at all places with the ordinary ideal class
group.  See J.S. Milne, \emph{Class Field Theory}, v4.03 (2020),
Theorem~0.4 and V, Example~3.9,
\url{https://www.jmilne.org/math/CourseNotes/CFT.pdf}.}
Consequently
$\operatorname{Sha}^1_S(M^D)$ injects into
$\operatorname{Hom}(\operatorname{Cl}(\mathcal O_L),\mathbb F_\ell)$.
Each cyclic factor of the class group contributes either $1$ or
$\ell$ to the latter order, never more than its own order.
Thus $|\operatorname{Sha}^1_S(M^D)|\le h_L$.
Combine this with (KB.15), $|S|\le n$, and the injection (KB.4).
This proves exactly the bound used by \texttt{card\_mod\_le},
with no switch to a larger number field or a different $D$.
\end{proof}

\subsection{The complete quantitative join}

\begin{proposition}[KB7: the literal reduction inequality]
\label{claim:KB7}
If $T$ of (KB.2) is finite, then $X$ is finite and
\[
 \log |X|\le
 4n\log(2\Delta)\,(2e)(2e+3)\log(n+1)+\log|T|.
 \tag{KB.16}
\]
All constants in this inequality are independent of
$\ell$, $\Omega$, and the local choices in $D$.
\end{proposition}

\begin{proof}
We include the elementary class-number estimate needed to control
the first factor in (KB.13).  Minkowski's theorem supplies an
integral ideal in each class of norm at most
$\frac{n!}{n^n}(4/\pi)^{r_2}\sqrt\Delta\le\sqrt\Delta$.
For the last inequality use $r_2\le n/2$ and, for $n\ge2$,
$n!\le((n+1)/2)^n$; the resulting factor is at most
$(3/(2\sqrt\pi))^n<1$.  For $n=1$ it equals one.\footnote{
Minkowski's ideal-class representative bound is stated with this
normalization in J.S. Milne, \emph{Algebraic Number Theory},
v3.08 (2020), Theorem~4.3,
\url{https://www.jmilne.org/math/CourseNotes/ANT.pdf}.}
At a rational prime the number of prime ideals is at most $n$ and
their residue degrees are positive.  Its ideal-counting Euler
factor is therefore coefficientwise bounded by $(1-t)^{-n}$.
Multiplication over rational primes shows that the number of
ideals of norm $a$ is at most the number $d_n(a)$ of ordered
factorizations $a=a_1\cdots a_n$ into positive integers.
For $B\ge1$, summing one variable at a time and using
$\sum_{a\le B}1/a\le1+\log B$ gives
\[
 \sum_{a\le B}d_n(a)\le B(1+\log B)^{n-1},\qquad
 \log h_L\le\frac12\log\Delta
 +(n-1)\log(1+\tfrac12\log\Delta)
 \le\frac n2\log\Delta.
 \tag{KB.17}
\]
Here $\Delta\ge1$ and $\log(1+t)\le t$ for $t\ge0$.

Put $a=|X/\ell X|$.  KB6 and (KB.17), with $\ell\ge2$, imply
\[
 \log a\le\log h_L+n\log\ell
                  \le4n\log(2\Delta)\log\ell.
 \tag{KB.18}
\]
For example, $\log\ell\ge\log2>1/2$ bounds
$\tfrac12\log\Delta$ by $\log\Delta\log\ell$, and
$4\log2>1$ supplies the remaining $\log\ell$ term.
Also $\ell\le n+1$.  For odd $\ell$, the primitive root gives
$\mathbb Q(\mu_\ell)\subset L$ and degree $\ell-1\le n$;
the degree formula follows from Eisenstein applied to
$\Phi_\ell(Y+1)$.  For $\ell=2$ use $n\ge1$.

KB4 and KB5 now give
$|X|\le a^{c(2e)}|T|$.  All group orders are positive, so
logarithms preserve the inequality.  Since
$c(m)=m(m+5)/2\le m(m+3)$ and
$\log\ell\le\log(n+1)$, (KB.18) gives (KB.16).
This is the numerical inequality in \texttt{Reduction.reduction},
including its homology rank $4e+4$ and cohomological degree $2e$.
\end{proof}

\begin{proposition}[KB8: one universal constant and the literal LemmaK endpoint]
\label{claim:KB8}
For any real $C$, \texttt{TorsBoundPow C} implies
\[
              \texttt{LemmaK}\bigl(\max(C,0)+9\bigr).
 \tag{KB.19}
\]
In particular TB10, which supplies \texttt{TorsBoundPow 22}, gives
$\texttt{LemmaK 31}$.  Explicitly, for every original arithmetic
datum in (KB.1),
\[
 X\text{ is finite},\qquad
 \log|X|\le
       \exp\!\left((2en)^{31}\log(2\Delta)\right).
 \tag{KB.20}
\]
Thus $\exists C_K>0\;\texttt{LemmaK}\ C_K$, with $C_K$ chosen
before $\ell,L,\Omega,D,e$.
\end{proposition}

\begin{proof}
Set $x=2en\ge2$, $A=\log(2\Delta)>0$, and
$p=x^{\max(C,0)}\ge1$.  The hypothesis gives finiteness of $T$
and $\log|T|\le\exp(x^C A)$.  The coefficient of $A$ in the
first term of (KB.16) satisfies
\[
 Q=4n(2e)(2e+3)\log(n+1)
 \le 4n(2e)(5e)n=10x^2\le x^9-1.
 \tag{KB.21}
\]
The first inequality uses $e\ge1$ and $\log(n+1)\le n$.
For the last, $x^7\ge2^7$ and $x^2\ge4$ give
$x^9-10x^2=x^2(x^7-10)\ge472>1$.
Therefore $QA\le B:=(x^9-1)pA$, while
$\exp(x^CA)\le\exp(pA)$.  Since $B\ge0$ and $\exp(pA)\ge1$,
\[
 \exp(pA)+B\le\exp(pA)(1+B)
                \le\exp(pA+B)
                =\exp\!\left(x^{\max(C,0)+9}A\right).
\]
Combining this with KB7 proves (KB.19).  The exponent is positive,
and it is a function only of the single constant $C$ in the
torsion theorem.  Substitution $C=22$ proves (KB.20) and gives the
literal quantified statement required downstream.
\end{proof}

\paragraph{Scope of the closure.}
KB4 is an actual supplier, rather than an implication whose
\texttt{KBridge} premise remains unfilled.  KB8 closes the whole
shared \texttt{LemmaK} branch once the separately proved TB10 is
joined: finite arithmetic $K$-theory, the correct continuous
cohomology, the finite-rank integral homology target, every kernel
exponent, the strict local conditions, and the universal constant
have all been connected.  The classical leaves just cited remain
citation leaves; no proof of those foundational theorems from
first principles is claimed.  This component asserts neither a
full RH/GRH reconstruction nor a replay of the Lean checker.
